% Generated by scripts/assemble.py; corrections belong in source/corrections.json.
\examhead{4}{北京理工大学 2025 年数学分析试卷(回忆版)}{11}
% p012-b08; source PDF page 12
\par\Needspace{5\baselineskip}
\blockanchor{p012-b08}
\majorquestion{一}计算题.\par
% p012-b09; source PDF page 12
\blockanchor{p012-b09}
\question{1}$\lim_{t\to 0}\frac{f(t) - 2}{t} = 3,$ 求 $\lim_{t\to 0}\frac{e^{f(t)} - 2}{t - 2}.$\par
% p012-b10; source PDF page 12
\blockanchor{p012-b10}
\question{2}$\lim_{x\to0}\left(\frac{1}{\sin^{2}x}-\frac{1}{x^{2}}\right)$ .\par
% p012-b11; source PDF page 12
\blockanchor{p012-b11}
\question{3}$f(x,y) = 7xy$ 在 $(x - 1)^2 +y^2 = 1$ 上的最大值.\par
% p012-b12; source PDF page 12
\blockanchor{p012-b12}
\question{4}$\iint_{D} (x + 2xy) \, \mathrm{d}x \, \mathrm{d}y$, 其中 $D: -1 \leq x \leq 1, -1 \leq y \leq x$.\par
% p012-b13; source PDF page 12
\blockanchor{p012-b13}
\question{5}$F(x) = \int_{1}^{x^2}\mathrm{d}t\int_{t^2}^{x^2}f(s,t)\mathrm{d}s,$ 求 $F^{\prime}(x),F^{\prime \prime}(x).$\par
% p012-b14; source PDF page 12
\blockanchor{p012-b14}
\question{6}$u = x^{2}y + y^{2}z + z^{2}x,$ 求 \displayformula{\frac{\partial^2u}{\partial x^2}\Big|_{(0,0,1)},\frac{\partial^2u}{\partial x\partial y}\Big|_{(1,0, - 1)},\frac{\partial^2u}{\partial x\partial z}\Big|_{(1, - 1, - 1)}.}\par
% p012-b15; source PDF page 12
\blockanchor{p012-b15}
\question{7}忘了.\par
% p012-b16; source PDF page 12
\par\Needspace{5\baselineskip}
\blockanchor{p012-b16}
\majorquestion{二}解答如下问题:\par
% p012-b17; source PDF page 12
\blockanchor{p012-b17}
\question{1}证明: $e^{-x} + e^{x} > x^{2} - 2, x \neq 0$.\par
% p012-b18; source PDF page 12
\blockanchor{p012-b18}
\question{2}$f(x)$ 连续, $\lim_{x\to 0}\frac{f(5x - x)}{x} = 3,$ 求 $f(x)$ 可导, 且 $f^{\prime}(0) = \frac{3}{4}$.\par
% p012-b19; source PDF page 12
\par\Needspace{5\baselineskip}
\blockanchor{p012-b19}
\majorquestion{三}求幂级数 $\sum_{n=0}^{\infty} \frac{x^{2n+1}}{2n+1}$ 的收敛半径, 收敛点集, 和函数.\par
% p012-b20; source PDF page 12
\par\Needspace{5\baselineskip}
\blockanchor{p012-b20}
\majorquestion{四}$u = f\left(\ln \sqrt{x^2 + y^2}\right), \frac{\partial^2 u}{\partial x^2} + \frac{\partial^2 u}{\partial y^2} = (x^2 + y^2)^{\frac{3}{2}}$，求 $f(t)$ 的表达式.\par
% p012-b21; source PDF page 12
\par\Needspace{5\baselineskip}
\blockanchor{p012-b21}
\majorquestion{五}$f(x) = \sin \frac{2}{x}$, 证明 $f(x)$ 在 $(a, +\infty)(a > 0)$ 上一致连续, 在 $(0, +\infty)$ 上不一致连续.\par
% p012-b22; source PDF page 12
\par\Needspace{5\baselineskip}
\blockanchor{p012-b22}
\majorquestion{六}$f(x)$ 非负二阶可导, 且存在 $M > 0$ 满足 $|f''(x)| \leq M$, 证明: $f'(x) \leq \sqrt{2Mf(x)}$.\par
% p012-b23; source PDF page 12
\par\Needspace{5\baselineskip}
\blockanchor{p012-b23}
\majorquestion{七}已知 $\lim_{x\to \infty}f(x) = f(0)$.\par
% p012-b24; source PDF page 12
\blockanchor{p012-b24}
(1) 证明: 存在 $\xi_1, \xi_2 \in (0, +\infty)$, $\xi_1 \neq \xi_2$, 使得 $f(\xi_1) = f(\xi_2)$.\par
% p012-b25; source PDF page 12
\blockanchor{p012-b25}
(2) $f(x)$ 至少有一个最大值或最小值.\par
% p012-b26; source PDF page 12
\blockanchor{p012-b26}
(3) $0 \leq f(x) \leq \frac{x^3}{e^x}$, 证明: 存在 $\xi$ 满足 $f'(\xi) = \frac{\xi^2(3 - \xi)}{e^\xi}$.\par
% p013-b00; source PDF page 13
\par\Needspace{5\baselineskip}
\blockanchor{p013-b00}
\majorquestion{八}$I(p) = \int_{0}^{+\infty}\frac{\cos^{2}x}{x^{p}}\mathrm{d}x.$\par
% p013-b01; source PDF page 13
\blockanchor{p013-b01}
(1) 求 $I(p)$ 的定义域.\par
% p013-b02; source PDF page 13
\blockanchor{p013-b02}
(2) 讨论 $I(p)$ 在定义域上的连续性.\par
% p013-b03; source PDF page 13
\blockanchor{p013-b03}
(3) 证明 $I(p)$ 在 $(-1, 1)$ 不一致收敛.\par
